\section{Introduction}
Direct-to-satellite IoT extends LoRa coverage to areas without terrestrial gateways by placing a
gateway on a low-Earth-orbit (LEO) satellite~\cite{fraire2019}. The link is only available during
passes of a few minutes, its SNR follows the elevation, and the carrier drifts by hundreds of hertz
per second near zenith. Adaptive data rate (ADR)~\cite{semtechADR} was designed for static
terrestrial links: it needs 20 uplinks (more than three hours at one reading every 10 minutes) and
never raises the spreading factor (SF) on its own. Events such as jamming, blockage, gateway outages
or the end of a satellite pass are outside what it was designed for.

Large language models (LLMs) that call tools~\cite{yao2023react} have been proposed for network
operation~\cite{maatouk2024}. This work asks a narrow question that can be measured:
\emph{can a small, local LLM agent, called only on events and acting through validated tools,
choose LoRa parameters better than ADR and a classifier trained on oracle labels when the rules do
not anticipate an event, and what does it cost in decision latency?}

\textbf{Contributions.}
(i) A packet-level model of LoRa under a LEO Doppler rate with two receiver classes, validated
against the exact noncoherent SER and compared with the 550\,km limit measured on a
CubeSat~\cite{doroshkin2019}.
(ii) An open hybrid terrestrial/LEO LoRa network simulator with a class-A control plane and common
random numbers across policies, so every comparison is paired.
(iii) An event-triggered LLM agent with five guarded tools, evaluated against fixed SFs, ADR, an ML
classifier, an oracle, and a rule-based controller that uses the same triggers and information.
(iv) A measurement of the agent's cost (latency, calls, invalid actions) on a laptop-class CPU.
All code, configurations, raw CSV files, agent traces and figure scripts are in the repository;
every number in this report is generated from them.
